\ncdmodultest{4}{Forbrug, Boligskift og Rejseoplevelser}{全真模考卷四：农场购物、搬入新居与难忘假期}{特色食品店与农场商店 / 莫勒一家搬进联排别墅 / 丹麦冬泳与桑拿文化 / 三人谈最难忘的假期 / 乔迁新居回信与向主管申请调班休假 / 实习与美食俱乐部口语及合租信息差}{全册五大语法支柱综合实战：分裂强调句 Det er... der/som/at · 动静三维方位副词 ud/ude/op/oppe/ind/inde · 礼貌过去时 kunne/ville/måtte · 直接引语转间接引语 · 复合句时态无缝切换}

% =====================================================================
% Del 1: Læsning (40 minutter · 24 point)
% =====================================================================
\ncdsubsec{Del 1: Læsning（阅读理解 · 40 分钟 · 24 分）}

\begin{tipbox}{考场规则与配速提醒（Læsning Opgave 1--4）}
Opgave 1, 2, 3 og 4 udleveres samlet. Der afsættes \textbf{40 minutter} i alt, og du må \textbf{ikke} bruge ordbog, noter eller mobiltelefon. Maksimum: \textbf{24 point}（及格线：\textbf{16--24 point = Bestået}；\textbf{0--15 point = Ikke bestået}）。建议时间分配：Opgave 1 (8 min) $\rightarrow$ Opgave 2 (8 min) $\rightarrow$ Opgave 3 (8 min) $\rightarrow$ Opgave 4 (12 min)，最后留 4 分钟核对涂卡与拼写。
\end{tipbox}

% ---------------------------------------------------------------------
\subsubsection*{Opgave 1 (NY · 6 point): \textit{Gårdbutikker og specialbutikker}}

\begin{emailbox}{Kursistark --- Instruktion（Opgave 1）}
\begin{itemize}
  \item Læs de seks små tekster om personerne, der gerne vil købe særlige madvarer eller gaver (1--6).
  \item Find den butik / annonce (B--I), der passer til hver tekst.
  \item Skriv bogstavet på annoncen, der passer til teksten. Der er \textbf{to annoncer, du ikke skal bruge}. Se eksemplet (0).
\end{itemize}

\vspace{0.3em}
\hrule
\vspace{0.4em}

\begin{enumerate}[label=\textbf{\arabic*.}, leftmargin=2.0em, itemsep=0.45em]
  \item[\textbf{0.}] Henrik og Mette vil gerne køre en tur ud på landet på lørdag. De vil købe friske, økologiske grøntsager og æg direkte fra gården, og så håber de også at kunne plukke nogle jordbær selv og drikke en kop varm kaffe. \hfill \textcolor{NCDRed}{\textbf{[Eksempel: Svar A]}}
  \item Sofie skal til sin søsters 40-års fødselsdag på fredag og leder efter en særlig gave. Søsteren elsker god kaffe og mørk chokolade, men hun drikker aldrig te. Sofie vil gerne have, at butikken pakker gaven pænt ind med et lille kort, så den er helt klar til festen. \hfill\mbox{\textbf{Svar: \rule{1.2cm}{0.4pt}}}
  \item Jonas og Amalie vil gerne bruge meget mindre plastik i hverdagen og smide mindre affald ud. Derfor leder de efter en butik, hvor de kan tage deres egne glas og stofposer med og kun betale for den mængde ris, nødder og krydderier, som de faktisk skal bruge. \hfill\mbox{\textbf{Svar: \rule{1.2cm}{0.4pt}}}
  \item Nadia spiser hverken kød, fisk eller mælkeprodukter. Hun har tit travlt efter arbejde på hospitalet og orker ikke at lave aftensmad fra bunden. Hun søger et sted i byen, hvor hun kan hente færdiglavede, plantebaserede retter og salater, som hun bare kan tage med hjem. \hfill\mbox{\textbf{Svar: \rule{1.2cm}{0.4pt}}}
  \item Søren har inviteret fire gamle venner til frokost på lørdag. Han vil gerne servere klassisk dansk frokost med frisk fisk direkte fra lokale fiskere. Han har ikke lyst til selv at stå i køkkenet hele morgenen, så han vil også kunne købe røget fisk og færdige fiskefrikadeller. \hfill\mbox{\textbf{Svar: \rule{1.2cm}{0.4pt}}}
  \item Freja og Kasper elsker at bruge søndag formiddag på en langsom morgenmad derhjemme. De leder efter et sted, der har åbent søndag morgen og sælger nybagt surdejsbrød samt specielle danske oste, som er lavet i hånden på små mejerier --- ikke almindelige varer fra et supermarked. \hfill\mbox{\textbf{Svar: \rule{1.2cm}{0.4pt}}}
  \item Birgitte skal ud at cykle en lang tur langs kysten på søndag og vil gerne gøre et stop på vejen. Hun har ikke brug for kaffe eller brød, men elsker naturprodukter som lokal honning til sine pandekager og friskpresset juice, og hun betaler helst med MobilePay. \hfill\mbox{\textbf{Svar: \rule{1.2cm}{0.4pt}}}
\end{enumerate}
\end{emailbox}

\vspace{0.3em}
{\small
\renewcommand{\arraystretch}{1.3}
\begin{longtable}{@{} >{\raggedright\arraybackslash}p{8.2cm} @{\hspace{0.6cm}} >{\raggedright\arraybackslash}p{8.2cm} @{}}
\toprule
\multicolumn{2}{@{} >{\raggedright\arraybackslash}p{16.95cm} @{}}{\textbf{\color{NCDNavy}Annoncer for gårdbutikker og specialbutikker (A--I)}} \\
\midrule
\endfirsthead
\toprule
\multicolumn{2}{@{} >{\raggedright\arraybackslash}p{16.95cm} @{}}{\textbf{\color{NCDNavy}Annoncer for gårdbutikker og specialbutikker (fortsat)}} \\
\midrule
\endhead
\bottomrule
\endfoot

\multicolumn{2}{@{} >{\raggedright\arraybackslash}p{16.95cm} @{}}{%
  \textbf{\color{NCDRed}A. Lindegaard Øko-Gårdbutik [Eksempel]}\newline
  Besøg vores hyggelige økologiske gårdbutik blot 15 minutter uden for byen! Vi sælger friskhøstede kartofler, gulerødder og æg fra vores fritgående høns. Fra juni til august kan hele familien gå ud på marken og selv plukke søde jordbær. I vores lille gårdcafé serverer vi nybrygget kaffe og hjemmebagt kage. Åbent tirsdag til søndag kl. 10--17. Se mere på \textit{lindegaard-oko.dk}%
} \\[6pt]
\midrule

\textbf{\color{NCDNavy}B. Kystens Fiskehus}\newline
Friskfanget fisk hver dag direkte fra de lokale fiskekuttere i havnen! Hos os kan du købe alt fra rødspætter og torsk til nygravet laks. Har du travlt, eller skal du have gæster til frokost? Prøv vores hjemmelavede fiskefrikadeller, varmrøget makrel eller rejecocktail, som er lige til at sætte på bordet. Vi har åbent tirsdag til lørdag kl. 9--17.30. Bestil eventuelt på forhånd på \textit{kystensfisk.dk}
&
\textbf{\color{NCDNavy}F. Kaffe \& Kakao Kompagniet}\newline
Vi rister vores egne specialkaffebønner fra små gårde i Sydamerika og laver håndlavet mørk og lys chokolade i vores eget værksted. Leder du efter en flot værtindegave eller fødselsdagsgave? Vi samler gerne en personlig gaveæske efter dit ønske og pakker den smukt ind med silkebånd og et håndskrevet hilsenkort. Kig forbi i gågaden mandag--lørdag, eller bestil på \textit{kaffekakao.dk} \\[6pt]
\midrule

\textbf{\color{NCDNavy}C. Møllebakkens Bageri \& Ost}\newline
Vi elsker godt håndværk! Hver morgen bager vi sprødt surdejsbrød og grove rundstykker i vores stenovn af økologisk mel. I vores store køledisk finder du desuden over 40 forskellige håndlavede oste fra små danske gårdmejerier samt friskkærnet smør. Kig ind efter morgenbrød og lækker ost --- også lørdag og søndag fra kl. 7.00! Du finder os på Møllevej 14 i midtbyen.
&
\textbf{\color{NCDNavy}G. Slagtermester Holm}\newline
Den klassiske danske slagter i hjertet af byen siden 1984! Vi tilbyder mørt okse- og svinekød fra danske landmænd, hjemmelavede grillpølser samt lækkert pålæg og lun leverpostej til familiens madpakker. Skal du holde sommerfest eller grillmiddag, laver vi gerne store kødpakker klar til grillen. Åbent tirsdag--fredag kl. 9--17.30 og lørdag kl. 8--14. Besøg os på \textit{slagterholm.dk} \\[6pt]
\midrule

\textbf{\color{NCDNavy}D. Grønt \& Glas --- Emballagefri Butik}\newline
Vil du spare på plastik og undgå madspild? I vores nul-affaldsbutik kan du medbringe dine egne rene glas, dåser og stofposer hjemmefra. Du vejer blot beholderen ved indgangen og fylder derefter præcis den mængde ris, pasta, linser, nødder, havregryn eller krydderier i, som du har brug for. Vi har også miljøvenlig sæbe på tap. Åbent mandag--lørdag kl. 10--18. Læs mere på \textit{groentogglas.dk}
&
\textbf{\color{NCDNavy}H. Sønderstrand Honning \& Æblemost}\newline
Velkommen til vores lille familieejede frugtplantage lige ud til den smukke kystcykelrute! I vores hyggelige gårdbutik sælger vi ren, lokal blomsterhonning fra egne bier, koldpresset æblemost af nyplukkede æbler og hjemmelavet bærmarmelade. Gårdbutikken har åbent hver søndag kl. 10--16, og resten af ugen kan du handle i vores døgnåbne vejbod med selvbetjening og nem betaling via MobilePay. \\[6pt]
\midrule

\textbf{\color{NCDNavy}E. Sydhavnens Grønne Delikatesse}\newline
100 \% plantebaseret mad lige til at tage med hjem! Hver dag laver vi friske måltidssalater, varme simreretter, cremet hummus og sprøde sandwich helt uden kød, æg eller komælk. Perfekt til dig, der ønsker en hurtig, sund og lækker aftensmad efter en lang arbejdsdag. Alle retter pakkes i bionedbrydelige bokse. Åbent alle hverdage kl. 11--19. Se ugens menu på \textit{sydgroen.dk}
&
\textbf{\color{NCDNavy}I. Nordisk Krydderi \& Te-hus}\newline
Elsker du en god kop te? Vi har mere end 90 slags løs te --- alt fra grøn japansk matcha og hvid te til krydret chai --- samt eksotiske krydderier fra Indien og Marokko. Alle vores teer og krydderier sælges færdigpakkede i flotte, lufttætte metaldåser, så duften bevares bedst muligt. Vi sælger også håndlavede tekander. Åbent tirsdag--lørdag kl. 10--17 på \textit{teogkrydderi.dk} \\
\end{longtable}
}

% ---------------------------------------------------------------------
\subsubsection*{Opgave 2 (5 point): \textit{Familien Møller flytter i rækkehus}}

\begin{emailbox}{Kursistark --- Instruktion（Opgave 2）}
Læs teksten \textit{Familien Møller flytter i rækkehus}. I hvert afsnit er der \textbf{én sætning, der ikke passer} i afsnittet. Find den sætning, der ikke passer. \textbf{Streg sætningen ud.} Se eksemplet i afsnit (0).

\vspace{0.35em}
\hrule
\vspace{0.4em}

\textbf{(0)} Efter fem år i en lille toværelses lejlighed på fjerde sal i Aarhus har Thomas og Katrine Møller endelig købt et lyst rækkehus i en rolig forstad. Deres to børn, Ida og Villads, glæder sig helt vildt til at få deres egne værelser og en lille have. \textcolor{NCDRed}{\textbf{[Eksempel --- streges ud:]}}\ \textit{De vil aldrig flytte væk fra fjerde sal.} Hele familien har pakket papkasser i tre uger og gjort klar til den store flyttedag.

\vspace{0.35em}
\textbf{(1)} Lørdag morgen klokken otte holdt en stor hvid flyttebil ude foran opgangen, og tre stærke venner kom for at hjælpe. Først bar de alle de tunge møbler og papkasser ned ad trapperne, selvom det var hårdt uden elevator. Thomas og Katrine bar alle tingene helt alene. Bagefter kørte de ud til det nye rækkehus, hvor de hurtigt satte alle kasserne ind i stuen.

\vspace{0.35em}
\textbf{(2)} Midt på formiddagen ringede det pludselig på døren i det nye hus. Ude på trappen stod de nye naboer, Jens og Helle, med en varm termokande kaffe og en kurv med friskbagte boller. Alle fortsatte med at bære kasser uden at kigge på naboerne. Familien Møller blev meget rørte over den søde velkomst og holdt straks en kort pause i haven. De snakkede hyggeligt om kvarteret i tyve minutter, før de gik ind igen.

\vspace{0.35em}
\textbf{(3)} Om eftermiddagen begyndte familien at pakke ud og sætte møblerne på plads. Katrine hang gardiner op i stuen, mens Thomas samlede børnenes senge oppe på første sal. Ida og Villads valgte straks, hvilket værelse der skulle være deres, og satte deres bøger på hylden. Børnene blev kede af det, fordi de stadig skulle dele et lille værelse. Klokken seks var de fleste møbler endelig stillet op.

\vspace{0.35em}
\textbf{(4)} Sidst på dagen var alle både sultne og helt udmattede efter det hårde arbejde. Det var umuligt at finde gryder og pander mellem de mange uåbnede kasser i køkkenet. Derfor lavede Katrine en stor og lækker middag med tre retter. I stedet ringede Thomas til det lokale pizzeria og bestilte fem varme pizzaer, som de hentede ti minutter senere. De spiste dem direkte af æskerne på gulvet i stuen.

\vspace{0.35em}
\textbf{(5)} Søndag formiddag gik familien en tur rundt i området for at lære deres nye kvarter bedre at kende. De opdagede, at der kun var fem minutters gang ned til en lille sø og en stor legeplads til børnene. Skolen lå desuden lige på den anden side af parken, så Ida kunne snart cykle derhen selv. Familien blev skuffet, fordi der slet ikke fandtes grønne områder eller legepladser i nærheden. Efter gåturen satte de sig tilfredse ud på deres nye terrasse i solen.
\end{emailbox}

% ---------------------------------------------------------------------
\subsubsection*{Opgave 3 (6 point): \textit{Vinterbadning og sauna --- varme i kulden}}

\begin{emailbox}{Kursistark --- Instruktion（Opgave 3）}
Læs teksten \textit{Vinterbadning og sauna --- varme i kulden}. Skriv de ord, der mangler. Du skal kun bruge hvert ord \textbf{én gang}. Der er \textbf{fire ord, du ikke skal bruge}. Se eksemplet i første linje.

\vspace{0.35em}
\hrule
\vspace{0.4em}

\textbf{Vinterbadning og sauna --- varme i kulden}

\vspace{0.25em}
Vinterbadning er i dag en meget \underline{\textbf{populær}} hobby i hele Danmark. Fra oktober til april går tusindvis af danskere ned til havet tidligt om morgenen, selvom vandet tit er meget \textbf{(1)}\ \rule{2.1cm}{0.4pt}.

\vspace{0.3em}
Først lægger man sit håndklæde og sin badekåbe klar i omklædningsrummet, \textbf{(2)}\ \rule{2.1cm}{0.4pt} man går forsigtigt ud ad træbroen. Det føles selvfølgelig iskoldt de første sekunder, når man hopper ned i bølgerne, \textbf{(3)}\ \rule{2.1cm}{0.4pt} bagefter oplever mange en dejlig varme og ny energi i hele kroppen.

\vspace{0.3em}
Det er faktisk ikke kun det kolde vand, \textbf{(4)}\ \rule{2.1cm}{0.4pt} gør aktiviteten så tiltrækkende. Efter den korte svømmetur løber de fleste hurtigt ind i en varm sauna, \textbf{(5)}\ \rule{2.1cm}{0.4pt} de sidder sammen, snakker og drikker te. Mange vinterbadeklubber har efterhånden lange ventelister, \textbf{(6)}\ \rule{2.1cm}{0.4pt} flere og flere søger både sundhed og fællesskab i de mørke måneder. På den måde bliver den kolde danske vinter lidt lettere at komme igennem.

\vspace{0.5em}
\begin{center}
\fbox{\parbox{0.92\linewidth}{\centering\vspace{3pt}
\textbf{fordi} \quad\textbar\quad \textbf{ude} \quad\textbar\quad \textbf{koldt} \quad\textbar\quad \textbf{at} \quad\textbar\quad \textbf{populær} \quad\textbar\quad \textbf{især} \quad\textbar\quad \textbf{hvor} \quad\textbar\quad \textbf{men} \quad\textbar\quad \textbf{ikke} \quad\textbar\quad \textbf{inden} \quad\textbar\quad \textbf{der}%
\vspace{3pt}}}
\end{center}
\end{emailbox}

% ---------------------------------------------------------------------
\subsubsection*{Opgave 4 (NY · 7 point): \textit{Tre personer fortæller om deres bedste ferieoplevelse}}

\begin{emailbox}{Kursistark --- Instruktion \& Svarskema（Opgave 4）}
\begin{itemize}
  \item Læs de tre tekster, og læs spørgsmålene (1--7).
  \item Find den person, der passer til hvert af de syv spørgsmål.
  \item Sæt kryds (\textbf{X}) ud for den person, der passer til spørgsmålet. Se eksemplet (0).
\end{itemize}

\vspace{0.3em}
{\small
\renewcommand{\arraystretch}{1.3}
\begin{tabularx}{\linewidth}{@{} >{\bfseries}p{0.6cm} >{\raggedright\arraybackslash}X >{\centering\arraybackslash}p{1.65cm} >{\centering\arraybackslash}p{1.65cm} >{\centering\arraybackslash}p{1.65cm} @{}}
\toprule
 & \textbf{Spørgsmål} & \textbf{Mads} & \textbf{Line} & \textbf{Ali} \\
\midrule
0. & Hvem var sammen med gamle venner på ferien? & A \underline{\textbf{X}} & B \underline{\hspace{0.45cm}} & C \underline{\hspace{0.45cm}} \\
\midrule
1. & Hvem fik hjælp af en lokal person, da noget gik i stykker? & A \underline{\hspace{0.45cm}} & B \underline{\hspace{0.45cm}} & C \underline{\hspace{0.45cm}} \\
2. & Hvem havde dårligt vejr, som ændrede ferieplanerne? & A \underline{\hspace{0.45cm}} & B \underline{\hspace{0.45cm}} & C \underline{\hspace{0.45cm}} \\
3. & Hvem var lidt nervøs før turen, fordi personen rejste alene? & A \underline{\hspace{0.45cm}} & B \underline{\hspace{0.45cm}} & C \underline{\hspace{0.45cm}} \\
4. & Hvem sparede penge ved ikke at spise på restaurant? & A \underline{\hspace{0.45cm}} & B \underline{\hspace{0.45cm}} & C \underline{\hspace{0.45cm}} \\
5. & Hvem var på cykelferie med sin familie? & A \underline{\hspace{0.45cm}} & B \underline{\hspace{0.45cm}} & C \underline{\hspace{0.45cm}} \\
6. & Hvem valgte sin rejseform for at passe på klimaet? & A \underline{\hspace{0.45cm}} & B \underline{\hspace{0.45cm}} & C \underline{\hspace{0.45cm}} \\
7. & Hvem talte med mange nye mennesker fra andre lande? & A \underline{\hspace{0.45cm}} & B \underline{\hspace{0.45cm}} & C \underline{\hspace{0.45cm}} \\
\bottomrule
\end{tabularx}
}
\end{emailbox}

\begin{emailbox}{Tre personer fortæller om deres bedste ferieoplevelse (A--C)}
\textbf{\color{NCDNavy}A. Mads}\\
Min bedste ferie var i oktober sidste år, da jeg lejede et gammelt sommerhus i Nordjylland sammen med fire gamle studiekammerater fra universitetet. Vores plan var at gå lange ture ved Vesterhavet og måske endda bade, men det regnede og blæste helt vildt næsten hver eneste dag. Derfor blev vi nødt til at ændre alle vores planer og blive indenfor. I stedet for at gå på stranden tændte vi op i brændeovnen og lavede store middage sammen. Feriehuset var billigt, og vi sparede mange penge, fordi vi købte ind i et lille supermarked og lavede al mad selv i stedet for at spise ude. Om aftenen sad vi ved bordet med varm kakao og spillede brætspil til langt ud på natten. Jeg har aldrig grinet så meget på en ferie før!

\vspace{0.45em}
\textbf{\color{NCDNavy}B. Line}\\
Sidste sommer rejste jeg helt alene med toget fra København op til fjeldene i Norge i ti dage. Jeg valgte toget i stedet for flyet, fordi jeg gerne ville rejse mere klimavenligt og nyde udsigten gennem vinduet. Før afrejsen var jeg faktisk lidt nervøs, fordi jeg aldrig havde været på ferie alene før. Jeg var bange for at føle mig ensom, men det skete heldigvis slet ikke! Vejret var klart og solrigt, og jeg vandrede femten kilometer om dagen gennem de grønne dale. Om aftenen sov jeg i forskellige små træhytter, hvor jeg mødte vandrere fra både Tyskland, Canada og Japan. Vi spiste varm suppe sammen og snakkede om vores rejser. Jeg kom hjem med ømme ben, men fyldt med selvtillid.

\vspace{0.45em}
\textbf{\color{NCDNavy}C. Ali}\\
I juli tog min kone, vores to børn på ni og tolv år og jeg på cykelferie til Bornholm. Vi lod bilen stå derhjemme og tog færgen fra Ystad med vores fire cykler. Hele ugen boede vi i et stort telt på den samme hyggelige campingplads tæt på Gudhjem. Solen skinnede fra morgen til aften, så vi cyklede omkring tyve kilometer hver dag for at se klipper, gamle kirker og små havne. Eftermiddagene brugte børnene på at bade i havet, mens vi spiste lækre fiskefrikadeller fra et lokalt røgeri. En dag punkterede min yngste søns cykel langt ude på en markvej, men en venlig landmand kørte forbi, hjalp os med at lappe hjulet og gav børnene friske kirsebær. Det glemmer vi aldrig!
\end{emailbox}

% =====================================================================
% Del 2: Skrivning (minimum 70 ord)
% =====================================================================
\clearpage
\ncdsubsec{Del 2: Skrivning（书面表达 · minimum 70 ord）}

\subsubsection*{Opgave 1（官方必考私人邮件回复）}

\begin{emailbox}{Kursistark --- Skrivning, Opgave 1}
\textbf{Situation:} Du har fået en mail fra din ven, Laura.\\
\textbf{Opgave:} Læs e-mailen, og skriv et svar til Laura. Du skal svare på alle Lauras spørgsmål. \textbf{Du skal skrive minimum 70 ord.}

\vspace{0.3em}
\hrule
\vspace{0.3em}

\textit{Hej\\
Jeg har hørt, at du endelig er flyttet! Stort tillykke! Hvordan ser din nye bolig ud? Hvordan var flyttedagen, og hjalp nogen dig? Hvad er du gladest for ved det nye kvarter? Og hvornår holder du indflytningsfest?\\
Mange hilsner\\
Laura}

\vspace{0.25em}
\hrule
\vspace{0.25em}
\textbf{Til:} Laura \qquad \textbf{Emne:} Ny bolig\\[2pt]
\rule{\linewidth}{0.4pt}\\[4pt]
\rule{\linewidth}{0.4pt}\\[4pt]
\rule{\linewidth}{0.4pt}\\[4pt]
\rule{\linewidth}{0.4pt}\\[4pt]
\rule{\linewidth}{0.4pt}\\[4pt]
\rule{\linewidth}{0.4pt}\\[4pt]
\rule{\linewidth}{0.4pt}
\end{emailbox}

\subsubsection*{Bonus Opgave 1B（半正式职场沟通邮件 · 礼貌过去时综合演练）}

\begin{emailbox}{Bonusopgave --- Skrivning til din chef om fridage}
\textbf{Situation:} Du arbejder på deltid i en boghandel, og din chef hedder Anne.\\
\textbf{Opgave:} Skriv en høflig e-mail til Anne (\textbf{minimum 70 ord}), hvor du:
\begin{itemize}[itemsep=0.12em, topsep=0.15em]
  \item bruger \textbf{høflig datid} (\textit{ville gerne spørge, om jeg måtte... / kunne det måske...}) til at bede om \textbf{3 dages fri næste måned (torsdag--lørdag)}, fordi du skal til en stor familiefest i udlandet;
  \item fortæller, at din kollega \textbf{Kasper} allerede har sagt ja til at tage din torsdagsvagt;
  \item tilbyder, at du \textbf{til gengæld} gerne vil tage et par ekstra vagter senere næste måned.
\end{itemize}

\vspace{0.25em}
\hrule
\vspace{0.25em}
\textbf{Til:} Anne (Chef) \qquad \textbf{Emne:} Ønske om fridage i næste måned\\[2pt]
\rule{\linewidth}{0.4pt}\\[4pt]
\rule{\linewidth}{0.4pt}\\[4pt]
\rule{\linewidth}{0.4pt}\\[4pt]
\rule{\linewidth}{0.4pt}
\end{emailbox}

% =====================================================================
% Del 3: Mundtlig kommunikation
% =====================================================================
\clearpage
\ncdsubsec{Del 3: Mundtlig kommunikation（口语交际全真题卡）}

\subsubsection*{Mundtlig Opgave 1: \textit{Emneopgave med Mindmap}}

\begin{tipbox}{官方考场规则提醒（Mundtlig Opgave 1）}
考生须自带 \textbf{2 张} 已填好的 Mindmap（其中至少 1 张属于工作 / 实习 / 学习 / 语言学校领域）。\textbf{铁律：Mindmap 上只能写关键词或短语（stikord），绝不能写完整句子！} 考场抽签选定 1 张：Del 1 脱稿陈述 \textbf{1--2 分钟}（考官看空白纸记笔记），Del 2 考官围绕你的内容追问 \textbf{3--4 分钟}。两部分均须及格。
\end{tipbox}

\begin{grammatikbox}{Topic 4A：Min praktik（官方 6 辐射格 · 纯 Stikord 备考卡）}
\small
\renewcommand{\arraystretch}{1.32}
\begin{tabularx}{\linewidth}{@{} >{\bfseries\color{NCDNavy}\raggedright\arraybackslash}p{4.1cm} >{\raggedright\arraybackslash}X @{}}
\toprule
\textbf{Mindmap 官方分支} & \textbf{Stikord 关键词笔记（禁止整句！）} \\
\midrule
1. praktiksted & lille it-firma i midtbyen · 12 ansatte · lyst kontor \\
2. arbejdsopgaver & teste nye hjemmesider · svare på kundemails · morgenmøde \\
3. dage / tid & tirsdag--torsdag · kl. 9.00--15.00 · 18 timer om ugen \\
4. kollegaer / chef & hjælpsomme kollegaer · venlig chef (Søren) · kun dansk på kontoret \\
5. pauser & frokost kl. 12.00 · 30 min. i kantinen · fælles kaffe \\
6. transport til praktiksted & cykel 20 minutter · bus i regnvejr \\
\bottomrule
\end{tabularx}

\vspace{0.45em}
\textbf{\color{NCDRed}Del 2 考官追问预测（Eksaminators opfølgende spørgsmål）：}
\begin{enumerate}
  \item \textit{Du siger, at I kun taler dansk på kontoret. Hvordan er det for dig, når dine kollegaer snakker hurtigt?}
  \item \textit{Hvilken arbejdsopgave synes du er mest spændende --- eller sværest --- på dit praktiksted?}
  \item \textit{Hvad spiser og snakker I om i frokostpausen i kantinen?}
  \item \textit{Kunne du tænke dig at arbejde fast i firmaet eller i en lignende virksomhed efter din praktik?}
\end{enumerate}
\end{grammatikbox}

\begin{grammatikbox}{Topic 4B：Min fritidsaktivitet: Madklub（官方 6 辐射格 · 纯 Stikord 备考卡）}
\small
\renewcommand{\arraystretch}{1.32}
\begin{tabularx}{\linewidth}{@{} >{\bfseries\color{NCDNavy}\raggedright\arraybackslash}p{4.1cm} >{\raggedright\arraybackslash}X @{}}
\toprule
\textbf{Mindmap 官方分支} & \textbf{Stikord 关键词笔记（禁止整句！）} \\
\midrule
1. sted / by & medborgerhus i Valby · stort fælleskøkken \\
2. dage / tid & hver anden onsdag · kl. 17.30--21.00 \\
3. pris / gratis & 80 kr. pr. gang · råvarer inklusiv \\
4. alene / sammen med & 10 faste medlemmer · nye og gamle venner · forskellige aldre \\
5. aktiviteter & lave mad fra hele verden · dække bord · spise sammen · rydde op \\
6. transport til aktiviteten & S-tog i 10 minutter · 5 min. gang fra stationen \\
\bottomrule
\end{tabularx}

\vspace{0.45em}
\textbf{\color{NCDRed}Del 2 考官追问预测（Eksaminators opfølgende spørgsmål）：}
\begin{enumerate}
  \item \textit{Du siger, at I laver mad fra hele verden. Hvad lavede I sidste gang, du var i madklubben?}
  \item \textit{Hvordan fandt du madklubben, og hvor længe har du været med?}
  \item \textit{Hvem bestemmer, hvad I skal lave --- og hjælper alle med at rydde op bagefter?}
  \item \textit{Hvad betyder det sociale fællesskab i klubben for dig og dit danske sprog?}
\end{enumerate}
\end{grammatikbox}

\clearpage
\subsubsection*{Mundtlig Opgave 2: \textit{Informationskløft --- Case: Mikkel}（满分 6 分 · 4--6 分及格）}

\begin{tipbox}{信息差规则提醒（Informationskløft）}
两人一组面对面抽取同一 Case，\textbf{绝对不能偷看对方卡片}。只有\textbf{提问方（Kursist A）在本题计分}：必须紧扣 2 个主题下共 6 个图片提示格，用第三人称 \textit{han / Mikkel} 提出 6 个聚焦的是非问句或特殊问句（\textbf{严禁对着右上角 Mikkel 头像发问！}）。每个合格提问得 1 分，共 6 分，\textbf{4--6 分及格}。随后双方互换角色抽新卡。
\end{tipbox}

\begin{emailbox}{Kursist A Ark --- Case: Mikkel（提问方提示卡，看不到 B 卡内容）}
\centering
\textbf{\color{NCDNavy}\large Case: Mikkel} \quad {\small\color{NCDGray}[Portrætfoto af Mikkel øverst --- stil IKKE spørgsmål til dette billede!]}

\vspace{0.4em}
\raggedright
\small
\renewcommand{\arraystretch}{1.35}
\begin{tabularx}{\linewidth}{@{} >{\raggedright\arraybackslash}X @{\hspace{0.5cm}} >{\raggedright\arraybackslash}X @{}}
\toprule
\textbf{\color{NCDNavy}Emne 1: Mikkel bor i et bofællesskab} & \textbf{\color{NCDNavy}Emne 2: Mikkels fritid / frivillige arbejde} \\
\midrule
\textbf{1. sted / type bolig} \hfill \textbf{\color{NCDRed}?}\newline
\textit{[Billede: et rækkehus på en grøn vej]} &
\textbf{4. frivilligt arbejde / sted} \hfill \textbf{\color{NCDRed}?}\newline
\textit{[Billede: en hyggelig lektiecafé på et bibliotek]} \\[4pt]
\midrule
\textbf{2. beboere / værelser} \hfill \textbf{\color{NCDRed}?}\newline
\textit{[Billede: en stue og et stort spisebord med flere stole]} &
\textbf{5. dage / tid} \hfill \textbf{\color{NCDRed}?}\newline
\textit{[Billede: en vægkalender og et ur]} \\[4pt]
\midrule
\textbf{3. regler / rengøring} \hfill \textbf{\color{NCDRed}?}\newline
\textit{[Billede: spand, kost og et ugeskema på væggen]} &
\textbf{6. andre hobbyer} \hfill \textbf{\color{NCDRed}?}\newline
\textit{[Billede: et par løbesko og en akustisk guitar]} \\
\bottomrule
\end{tabularx}
\end{emailbox}

\begin{emailbox}{Kursist B Ark --- Case: Mikkel（回答方图片信息卡，仅 B 可见）}
\small
\renewcommand{\arraystretch}{1.35}
\begin{tabularx}{\linewidth}{@{} >{\raggedright\arraybackslash}X @{\hspace{0.5cm}} >{\raggedright\arraybackslash}X @{}}
\toprule
\textbf{\color{NCDTeal}Emne 1: Mikkel bor i et bofællesskab} & \textbf{\color{NCDTeal}Emne 2: Mikkels fritid / frivillige arbejde} \\
\midrule
\textbf{1.} \textit{[Foto: et stort gult rækkehus i Roskilde, 10 min. fra stationen]}\newline
\textbf{Fakta:} Han bor i et stort rækkehus i Roskilde tæt på stationen. &
\textbf{4.} \textit{[Foto: skilt ”Røde Kors Lektiecafé --- Biblioteket”]}\newline
\textbf{Fakta:} Han arbejder frivilligt i en lektiecafé på biblioteket og hjælper børn med lektier. \\[4pt]
\midrule
\textbf{2.} \textit{[Foto: 5 unge voksne omkring et køkkenbord; navne på 5 døre]}\newline
\textbf{Fakta:} De bor 5 personer sammen; alle har eget værelse og deler køkken, stue og have. &
\textbf{5.} \textit{[Foto: opslag ”Mandag kl. 16.00--18.00”]}\newline
\textbf{Fakta:} Han er i lektiecaféen hver mandag eftermiddag fra kl. 16.00 til 18.00. \\[4pt]
\midrule
\textbf{3.} \textit{[Foto: tavle ”Fælles mad tirsdag \& torsdag” + ”Rengøring hver søndag”]}\newline
\textbf{Fakta:} De laver mad sammen to gange om ugen og gør rent hver søndag efter en fast plan. &
\textbf{6.} \textit{[Foto: en mand der løber i en skov og spiller guitar]}\newline
\textbf{Fakta:} Derudover løber han i skoven to gange om ugen og spiller guitar i et lille band. \\
\bottomrule
\end{tabularx}
\end{emailbox}

% =====================================================================
% Del 4: Retteark, Detaljeret Analyse & Højscore-Modeller
% =====================================================================
\clearpage
\ncdsubsec{Del 4: Retteark, Detaljeret Analyse \& Højscore-Modeller}

\subsubsection*{4.1 官方 \textit{Retteark} 评分卡与考前 30 秒终极自查清单}

\begin{grammatikbox}{DU3 Modul 2 --- Modultest 4 官方标准评分卡（Retteark）}
\small
\renewcommand{\arraystretch}{1.24}
\begin{tabularx}{\linewidth}{@{} >{\bfseries\color{NCDNavy}\raggedright\arraybackslash}p{3.45cm} >{\raggedright\arraybackslash}X >{\centering\arraybackslash}p{1.65cm} >{\centering\arraybackslash}p{2.55cm} @{}}
\toprule
\textbf{考核题型} & \textbf{标准答案（Rettenøgle）} & \textbf{满分} & \textbf{及格标准} \\
\midrule
Læsning Opgave 1 & \textbf{1.} F \quad \textbf{2.} D \quad \textbf{3.} E \quad \textbf{4.} B \quad \textbf{5.} C \quad \textbf{6.} H \quad（干扰项：G, I） & 6 point & 分项计分 \\
\cmidrule{1-4}
Læsning Opgave 2 &
\textbf{1.} \textit{Thomas og Katrine bar alle tingene helt alene.}\newline
\textbf{2.} \textit{Alle fortsatte med at bære kasser uden at kigge på naboerne.}\newline
\textbf{3.} \textit{Børnene blev kede af det, fordi de stadig skulle dele et lille værelse.}\newline
\textbf{4.} \textit{Derfor lavede Katrine en stor og lækker middag med tre retter.}\newline
\textbf{5.} \textit{Familien blev skuffet, fordi der slet ikke fandtes grønne områder eller legepladser i nærheden.} & 5 point & 分项计分 \\
\cmidrule{1-4}
Læsning Opgave 3 & \textbf{1.} koldt \quad \textbf{2.} inden \quad \textbf{3.} men \quad \textbf{4.} der \quad \textbf{5.} hvor \quad \textbf{6.} fordi \quad（干扰词：at, ikke, ude, især） & 6 point & 分项计分 \\
\cmidrule{1-4}
Læsning Opgave 4 & \textbf{1.} C \quad \textbf{2.} A \quad \textbf{3.} B \quad \textbf{4.} A \quad \textbf{5.} C \quad \textbf{6.} B \quad \textbf{7.} B & 7 point & 分项计分 \\
\cmidrule{1-4}
\textbf{\color{NCDRed}Læsning 总计} & \textbf{Opgave 1--4 合订总评}（40 分钟 · 禁用词典） & \textbf{24 point} & \textbf{\color{NCDRed}16--24 = Bestået} \\
\midrule
Skrivning Opgave 1 & 内容完整回答 Laura 全部 4 个问题 + 词数 $\ge$ 70 词 + 语序与基本语法清晰得体 & 合格制 & 建议目标 85--115 词 \\
\midrule
Mundtlig Opgave 1 & Del 1 脱稿陈述 1--2 分钟（仅带 stikord）+ Del 2 追问应答 3--4 分钟 & 合格制 & Del 1 \& 2 均须通过 \\
\midrule
Mundtlig Opgave 2 & 围绕 \textit{Mikkel} 的 6 个格子发出 6 个第三人称聚焦问题（不问头像） & 6 point & \textbf{4--6 = Bestået} \\
\bottomrule
\end{tabularx}
\end{grammatikbox}

\begin{tipbox}{考前 30 秒终极自查清单（DU3 Modul 2 结业实战总检）}
作为本册四套全真模考的收官卷，进考场前请默念以下 5 条保命检查项：
\begin{enumerate}
  \item \textbf{阅读同义转述铁律}：Opgave 1 和 Opgave 4 几乎永远用“换词表达（\textit{Parafrasering}）”设题——如“不喝牛奶” $\rightarrow$ \textit{uden komælk}；“关心气候” $\rightarrow$ \textit{rejse klimavenligt}；遇到原词复现要先警惕是不是干扰项！
  \item \textbf{写作 70+ 词四步验算}：\ding{172} 所有 4 个问句是否都已正面作答；\ding{173} 句首有时间/地点状语或前置从句时，主句是否执行 \textbf{V2 倒装}；\ding{174} 所有从句（\textit{fordi, at, som, når, da, om}）里的 \textit{ikke/altid/stadig} 是否放在\textbf{动词前}；\ding{175} 已结束过去时间用一般过去时，持续至今用 \textit{har/er} 完成时。
  \item \textbf{五大语法支柱绝不混用}：强调主语用 \textit{Det er..., der/som}，强调时间/地点用 \textit{Det er..., (at)}；间接是非问句用 \dabold{om}（不用 \textit{hvis}）；主语为复数（如 \textit{de}）禁止用 \textit{sin/sit/sine}。
  \item \textbf{动静方位与职场礼貌}：移动带去向用短形（\textit{ud til, ind i, op på}），静止所在用长形（\textit{ude ved, inde i, oppe på}）；请假/换班优先用礼貌过去时 \textit{ville/kunne/måtte}。
  \item \textbf{口试两条红线}：Opgave 1 只带纯 \textit{stikord}（不写整句）；Opgave 2 必须对准 6 个小格用第三人称 \textit{han/hun} 提问，切忌问右上角人物头像！
\end{enumerate}
\end{tipbox}

\subsubsection*{4.2 阅读 Opgave 1--4 逐题解析、干扰项拆解、全文中文精译与核心词表}

\begin{answerbox}
\textbf{\color{NCDNavy}【Opgave 1 逐题同义转述定位与干扰项排雷】}
\begin{itemize}
  \item \textbf{1. Svar: F (Kaffe \& Kakao Kompagniet)}：Sofie 需要“咖啡 + 黑巧克力 + 包装精美附带小卡片”，且明确表示姐姐“\textit{drikker aldrig te}”。广告 F 精准对应 \textit{specialkaffebønner}, \textit{håndlavet mørk... chokolade} 与 \textit{pakker den smukt ind med silkebånd og et håndskrevet hilsenkort}。\\
  \textbf{排雷}：干扰项 \textbf{I}（\textit{Nordisk Krydderi \& Te-hus}）主打 90 多种茶叶，恰好触犯“从不喝茶”的排除条件。
  \item \textbf{2. Svar: D (Grønt \& Glas --- Emballagefri Butik)}：Jonas og Amalie 想减少塑料垃圾（\textit{bruge meget mindre plastik / smide mindre affald ud}），自带玻璃罐和布袋（\textit{egne glas og stofposer}）按需称重购买大米、坚果和香料。广告 D 的 \textit{nul-affaldsbutik}, \textit{medbringe dine egne rene glas, dåser og stofposer}, \textit{vejer... præcis den mængde} 完美同义转述。\\
  \textbf{排雷}：干扰项 \textbf{I} 虽然也卖香料，但明确写明 \textit{sælges færdigpakkede i flotte, lufttætte metaldåser}（预包装金属罐），不能自带容器散装购买。
  \item \textbf{3. Svar: E (Sydhavnens Grønne Delikatesse)}：Nadia 既不吃肉、鱼也不吃奶制品（纯素），下班后没精力自己做饭，想买现成的植物基熟食带回家（\textit{færdiglavede, plantebaserede retter... tage med hjem}）。广告 E 对应 \textit{100 \% plantebaseret mad lige til at tage med hjem} 且明确 \textit{helt uden kød, fisk, æg eller komælk}。\\
  \textbf{排雷}：\textbf{C} 卖大量乳制品奶酪，\textbf{B} 卖鱼，\textbf{D} 卖的是需要回家自己煮的生米面干货。
  \item \textbf{4. Svar: B (Kystens Fiskehus)}：Søren 想请朋友吃经典丹麦午餐，需要本地渔船直供鲜鱼（\textit{direkte fra lokale fiskere}）以及熏鱼和现成鱼饼（\textit{røget fisk og færdige fiskefrikadeller}）。广告 B 精准命中 \textit{direkte fra de lokale fiskekuttere}, \textit{varmrøget makrel} 及 \textit{hjemmelavede fiskefrikadeller}，且周六营业。\\
  \textbf{排雷}：干扰项 \textbf{G}（\textit{Slagtermester Holm}）是传统肉铺，只卖牛猪肉、烤肠和猪肝酱，不卖鱼。
  \item \textbf{5. Svar: C (Møllebakkens Bageri \& Ost)}：Freja og Kasper 周日早上想买现烤酸面包（\textit{nybagt surdejsbrød}）以及小乳制品厂手工制作的非超市普通奶酪（\textit{specielle danske oste, lavet i hånden på små mejerier}）。广告 C 完美匹配 \textit{sprødt surdejsbrød}, \textit{håndlavede oste fra små danske gårdmejerier} 且 \textit{også lørdag og søndag fra kl. 7.00}。\\
  \textbf{排雷}：\textbf{H} 虽然周日营业，但只卖蜂蜜、苹果汁和果酱，没有面包和奶酪。
  \item \textbf{6. Svar: H (Sønderstrand Honning \& Æblemost)}：Birgitte 周日沿着海岸骑行（\textit{cykle... langs kysten på søndag}），不需要咖啡或面包，喜欢买本地蜂蜜和鲜榨果汁，最好能用 MobilePay 付款。广告 H 位于 \textit{kystcykelrute}，周日 10--16 点营业，出售 \textit{lokal blomsterhonning} 和 \textit{koldpresset æblemost}，并支持 \textit{MobilePay}。
\end{itemize}

\vspace{0.3em}
\textbf{\color{NCDTeal}[Opgave 1 中文精译]}\\
\textbf{人物需求 (0--6)}：(0) 亨里克和梅特周六想开车去乡下买新鲜有机蔬菜、鸡蛋，摘草莓并喝杯咖啡（对应 A）。(1) 索菲周五要参加姐姐 40 岁生日，想买好咖啡和黑巧克力的精美包装礼品附手写卡片（她姐姐从不喝茶）。(2) 乔纳斯和阿玛莉想减少塑料垃圾，找一家能自带玻璃罐和布袋按需称重买米、坚果和香料的店。(3) 娜迪亚不吃肉、鱼和奶制品，医院下班后没精力做饭，想买现成的植物基菜肴和沙拉带回家。(4) 瑟伦周六请四位老友吃经典丹麦午餐，想买本地渔民直供鲜鱼，以及现成的熏鱼和鱼饼。(5) 芙蕾雅和卡斯珀周日上午想吃慢悠悠的早饭，找周日早上开门、卖现烤酸面包和小农庄手工奶酪的店。(6) 比吉特周日沿海岸长途骑行，不需要咖啡或面包，想买本地蜂蜜配煎饼和鲜榨果汁，最好用 MobilePay 支付。\\
\textbf{店铺广告 (A--I)}：\textbf{A} 林德高有机农场商店（鲜蔬、散养鸡蛋、自摘草莓、农场咖啡）。\textbf{B} 海岸鱼屋（本地渔船直供鲜鱼、热熏鲭鱼、自制鱼饼，周二至周六）。\textbf{C} 磨坊坡面包与奶酪坊（石炉有机酸面包、40 多种小农庄手工奶酪，周六日早 7 点开门）。\textbf{D} 绿意与玻璃无包装商店（零浪费，自带容器称重买米面坚果香料）。\textbf{E} 南港绿色熟食店（100\% 植物基熟食、沙拉、鹰嘴豆泥外带，不含肉鱼蛋奶）。\textbf{F} 咖啡与可可公司（自烘咖啡豆、手工巧克力，定制丝带礼盒与手写贺卡）。\textbf{G} 霍尔姆肉铺（牛猪肉、烤肠、熟肉切片与猪肝酱，干扰项）。\textbf{H} 南滩蜂蜜与苹果汁果园（海岸骑行道旁，纯花蜜、冷榨苹果汁，周日营业及 MobilePay 自助摊）。\textbf{I} 北欧香料与茶屋（90 多种散装茶与金属罐装香料，干扰项）。

\vspace{0.45em}
\hrule
\vspace{0.45em}
\textbf{\color{NCDNavy}【Opgave 2 逐段矛盾句拆解与中文精译】}
\begin{itemize}
  \item \textbf{(1) 划掉：}\textit{Thomas og Katrine bar alle tingene helt alene.}——前句刚说三位强壮的朋友来帮忙，且用 \textit{de}（他们）一起把家具搬下楼，因此“两人独自搬完所有东西”完全矛盾。
  \item \textbf{(2) 划掉：}\textit{Alle fortsatte med at bære kasser uden at kigge på naboerne.}——后句明确写莫勒一家对新邻居的欢迎非常感动，立刻在花园小憩并愉快聊了二十分钟，故“大家看都不看邻居继续搬箱子”矛盾。
  \item \textbf{(3) 划掉：}\textit{Børnene blev kede af det, fordi de stadig skulle dele et lille værelse.}——第 (0) 段已交待孩子们对拥有“自己的房间”兴奋不已，本段前句也写 Ida 和 Villads 立刻挑好了各自房间，故“还要挤一间小房而难过”矛盾。
  \item \textbf{(4) 划掉：}\textit{Derfor lavede Katrine en stor og lækker middag med tre retter.}——前句说根本找不到锅碗瓢盆，后句也写 \textit{I stedet}（取而代之）去比萨店买了五张热比萨坐在客厅地板上吃，故“卡特琳娜做了一顿丰盛的三道菜晚餐”矛盾。
  \item \textbf{(5) 划掉：}\textit{Familien blev skuffet, fordi der slet ikke fandtes grønne områder eller legepladser i nærheden.}——前文刚写步行五分钟就有小湖、大游乐场和公园，后句也写一家人满意地坐在阳光露台上，故“因附近没有绿地游乐场而失望”矛盾。
\end{itemize}

\textbf{\color{NCDTeal}[Opgave 2 中文精译（已剔除干扰句）]}\\
(0) 在奥胡斯四楼一间狭小的两居室公寓住了五年后，托马斯和卡特琳娜·莫勒终于在安静的郊区买了一栋明亮的联排别墅。他们的两个孩子伊达和维拉斯兴奋极了，终于能拥有自己的房间和一座小花园。全家已经打包纸箱三个星期，为大搬家日做好了准备。(1) 周六早上八点，一辆白色大搬家卡车停在楼道门前，三位壮实的朋友前来帮忙。虽然没有电梯很辛苦，他们还是先把所有沉重的家具和纸箱搬下楼梯。随后他们开车前往新联排别墅，迅速把所有箱子搬进客厅。(2) 上午十点左右，新房子的门铃突然响了。门外台阶上站着新邻居延斯和赫勒，手里拿着一壶热咖啡和一篮刚出炉的面包卷。莫勒一家被这暖心的欢迎深深打动，立刻在花园里小歇片刻。大家愉快地聊了二十分钟关于这个街区的事，才又进屋。(3) 下午，全家人开始拆箱并把家具摆放到位。卡特琳娜在客厅挂上窗帘，托马斯则在二楼组装孩子们的床。伊达和维拉斯马上挑好了各自的房间，并把书摆上书架。到六点钟，大部分家具终于都摆放好了。(4) 一天结束时，大家经过辛苦劳作又饿又累。厨房里堆满了没拆封的纸箱，根本找不到锅和平底锅。于是托马斯打电话给当地比萨店订了五张热比萨，十分钟后去取了回来。大家坐在客厅地板上，直接从纸盒里拿比萨吃。(5) 周日上午，一家人在附近散步，想更熟悉自己的新街区。他们发现步行五分钟就能到一个小湖和一座大儿童游乐场。此外学校就在公园对面，伊达很快就能自己骑车去上学了。散步结束后，大家心满意足地坐到新露台上晒太阳。

\vspace{0.45em}
\hrule
\vspace{0.45em}
\textbf{\color{NCDNavy}【Opgave 3 逐空语法/语境解析与中文精译】}
\begin{itemize}
  \item \textbf{(1) koldt}：前文说从十月到四月冬泳，尽管海水经常非常“冷”（\textit{vandet} 为 \textit{et}-词，故形容词用中性 \textit{koldt}）。
  \item \textbf{(2) inden}：连接时间先后动作，“先在更衣室放好浴巾和浴袍，\textbf{在……之前}（\textit{inden}）小心走上木栈桥”。
  \item \textbf{(3) men}：语义转折，“最初几秒当然感觉冰冷刺骨，\textbf{但}（\textit{men}）之后很多人会感到浑身暖意与全新活力”。
  \item \textbf{(4) der}：分裂强调句（支柱一）！强调主语 \textit{det kolde vand}，后面直接接动词 \textit{gør}，因此只能用关系主语标记 \textbf{der}（绝不可用 \textit{at}！）。
  \item \textbf{(5) hvor}：地点定语从句，“跑进一间温暖的桑拿房，\textbf{在那里}（\textit{hvor}）大家一起坐着聊天喝茶”。
  \item \textbf{(6) fordi}：因果从句连词，解释为什么很多冬泳俱乐部有长长的候补名单，且后接从句语序 \textit{flere og flere søger...}，故填 \textbf{fordi}。
  \item \textbf{4 个干扰词排雷}：\textit{at}（不可引导主语型分裂句，亦不能充当第 5 空地点副词）；\textit{ikke}（副词无法连接两套主谓结构）；\textit{ude}（静态地点副词，第 1 空“\textit{meget ude}”与水温和 \textit{selvom} 均不搭配）；\textit{især}（副词，无法充当任何连词或形容词空位）。
\end{itemize}

\textbf{\color{NCDTeal}[Opgave 3 中文精译]}\\
冬泳如今在全丹麦是一项非常受欢迎的爱好。从十月到四月，成千上万的丹麦人一大早就会下到海边，尽管海水常常非常冰冷。首先，人们会在更衣室里把毛巾和浴袍放好，然后再小心翼翼地走上木栈桥。跳进海浪里的最初几秒当然感觉冰冷刺骨，但游完之后，许多人会体验到全身上下惬意的暖流和充沛的精力。实际上，让这项活动如此吸引人的并不仅仅是冰冷的海水。短时间游完泳后，大多数人会迅速跑进温暖的桑拿房，在那里一起坐着聊天喝茶。许多冬泳俱乐部渐渐排起了长长的候补名单，因为越来越多的人在阴暗的月份里寻求健康与社群归属感。这样一来，寒冷的丹麦冬天也变得更好熬一些了。

\vspace{0.45em}
\hrule
\vspace{0.45em}
\textbf{\color{NCDNavy}【Opgave 4 逐题同义定位与三人全文中文精译】}
\begin{itemize}
  \item \textbf{1. Svar: C (Ali)}：线索 \textit{da noget gik i stykker / hjælp af en lokal person} $\leftrightarrow$ Ali 儿子自行车爆胎（\textit{punkterede}），一位友善的当地农夫（\textit{en venlig landmand}）帮忙补胎。A 和 B 均未提及坏车或当地人施援。
  \item \textbf{2. Svar: A (Mads)}：线索 \textit{dårligt vejr, som ændrede ferieplanerne} $\leftrightarrow$ 原计划海滩散步游泳，但每天狂风暴雨（\textit{regnede og blæste helt vildt}），不得不改变计划待在室内（\textit{nødt til at ændre alle vores planer}）。B 和 C 全程晴朗。
  \item \textbf{3. Svar: B (Line)}：线索 \textit{lidt nervøs før turen... rejste alene} $\leftrightarrow$ Line 原文 \textit{Før afrejsen var jeg faktisk lidt nervøs, fordi jeg aldrig havde været på ferie alene før}。A 与同学同行，C 与妻儿同行。
  \item \textbf{4. Svar: A (Mads)}：线索 \textit{sparede penge ved ikke at spise på restaurant} $\leftrightarrow$ Mads 原文 \textit{vi sparede mange penge, fordi vi... lavede al mad selv i stedet for at spise ude}（\textit{spise ude} = 下馆子）。B 和 C 均未谈及省钱或不去餐馆。
  \item \textbf{5. Svar: C (Ali)}：线索 \textit{på cykelferie med sin familie} $\leftrightarrow$ Ali 带妻子和 9 岁、12 岁两个孩子骑车去博恩霍尔姆岛（\textit{min kone, vores to børn... på cykelferie}）。
  \item \textbf{6. Svar: B (Line)}：线索 \textit{valgte sin rejseform for at passe på klimaet} $\leftrightarrow$ Line 放弃飞机选择火车，是为了更环保低碳（\textit{valgte toget i stedet for flyet, fordi jeg gerne ville rejse mere klimavenligt}）。
  \item \textbf{7. Svar: B (Line)}：线索 \textit{talte med mange nye mennesker fra andre lande} $\leftrightarrow$ Line 在木屋里结识来自德国、加拿大和日本的徒步者并共进晚餐聊天（\textit{mødte vandrere fra både Tyskland, Canada og Japan... snakkede}）。Mads 只与老同学一起，Ali 只遇到一位当地农夫。
\end{itemize}

\textbf{\color{NCDTeal}[Opgave 4 中文精译]}\\
\textbf{A. 马兹}：我最难忘的一次假期是在去年十月，当时我和四位大学老同学一起在北日德兰租了一栋老度假屋。我们原本计划沿着北海长距离散步，甚至还想下海游泳，但几乎每天都在狂风暴雨。因此我们不得不改变所有计划，待在屋里。我们没去海滩，而是点起柴火炉，一起做大餐。度假屋很便宜，而且我们去小超市采购自己做所有饭菜而不下馆子，省了不少钱。晚上大家围坐在桌边喝热可可、玩桌游，一直玩到深夜。我从来没在度假时笑得这么开心过！\\
\textbf{B. 琳内}：去年夏天，我独自一人坐火车从哥本哈根北上挪威群山玩了十天。我选火车而不坐飞机，是因为我想以更环保的方式旅行，并透过车窗欣赏风景。出发前我其实有点紧张，因为我以前从没独自旅行过。我怕自己会感到孤单，但幸好完全没有！天气晴朗阳光明媚，我每天穿过翠绿的山谷徒步十五公里。晚上我住在不同的小木屋里，在那里遇到了来自德国、加拿大和日本的徒步爱好者。我们一起喝热汤，聊各自的旅程。回家时虽然腿酸，但我内心充满了自信。\\
\textbf{C. 阿里}：七月份，我和妻子以及九岁、十二岁的两个孩子一起去博恩霍尔姆岛骑行度假。我们把汽车留在家里，带着四辆自行车从于斯塔德乘渡轮登岛。整整一周我们都住在古德耶姆附近同一座温馨露营地的大帐篷里。太阳从早晒到晚，所以我们每天骑行约二十公里去看悬崖、古老教堂和小港口。下午孩子们去海里游泳，我们则吃当地熏鱼坊美味的鱼饼。有一天小儿子的自行车在偏僻乡间土路上爆胎了，一位友善的农夫刚好开车路过，帮我们补好车胎，还送给孩子们新鲜樱桃。那一幕我们永远不会忘记！
\end{answerbox}

\subsubsection*{本卷核心生词精表（Gloser --- Modultest 4）}

\begin{glosertable}
名词 & en gårdbutik & \textit{(gårdbutikken, gårdbutikker, -erne)} & 农场商店 \\
形容词 & emballagefri & \textit{(emballagefrit, emballagefrie)} & 无包装的，零浪费散装的 \\
名词 & et rækkehus & \textit{(rækkehuset, rækkehuse, -ene)} & 联排别墅 \\
名词 & en flyttedag & \textit{(flyttedagen, flyttedage, -ene)} & 搬家日 \\
动词 & at pakke ud & \textit{(pakker ud, pakkede ud, har pakket ud)} & 拆箱，打开包裹/行李 \\
名词 & vinterbadning & \textit{(vinterbadningen, ingen flertal)} & 冬泳 \\
名词 & et fællesskab & \textit{(fællesskabet, fællesskaber, -erne)} & 社群归属感，集体陪伴 \\
名词 & en brændeovn & \textit{(brændeovnen, brændeovne, -ene)} & 柴火炉 \\
名词 & en venteliste & \textit{(ventelisten, ventelister, -erne)} & 候补名单，排队名单 \\
动词 & at punktere & \textit{(punkterer, punkterede, er/har punkteret)} & 爆胎，扎破车胎 \\
形容词 & klimavenlig & \textit{(klimavenligt, klimavenlige)} & 气候友好的，低碳环保的 \\
名词 & en indflytningsfest & \textit{(indflytningsfesten, -fester, -erne)} & 乔迁派对，暖房聚会 \\
名词 & en vagt & \textit{(vagten, vagter, vagterne)} & 班次，值班 \\
动词 & at bytte & \textit{(bytter, byttede, har byttet)} & 交换，调换（班次） \\
名词 & et bofællesskab & \textit{(bofællesskabet, bofællesskaber, -erne)} & 合租共居之家 \\
\end{glosertable}

\subsubsection*{4.3 写作范文（A2 稳过版 + A2+ 高分示范版 + Bonus 1B 职场邮件）}

\begin{answerbox}
\textbf{\color{NCDNavy}【Opgave 1 范文一：A2 稳过版（85 ord · 清晰准确全覆盖）】}\\
\textit{Til: Laura \quad|\quad Emne: Ny bolig}

\vspace{0.25em}
\textit{Hej Laura}

\vspace{0.2em}
\textit{Tak for din mail! Ja, jeg er endelig flyttet ind i min nye lejlighed i Valby. Den har to lyse værelser, et lille køkken og en dejlig altan.}

\vspace{0.2em}
\textit{Flyttedagen i lørdags var hård, fordi vi bar alle møblerne op på tredje sal. Heldigvis hjalp min bror og to venner mig.}

\vspace{0.2em}
\textit{Jeg er gladest for, at kvarteret er så roligt, og at der ligger en park tæt på. Jeg holder indflytningsfest lørdag den 20. maj kl. 18. Jeg håber, du kommer!}

\vspace{0.2em}
\textit{Mange hilsner\\
Emma}

\vspace{0.25em}
\small\textbf{点评（85 词）}：句式简洁自然、零语法硬伤，完整回应了 Laura 的全部 4 个问题（房型外观 / 搬家日与谁帮忙 / 新街区最满意之处 / 乔迁派对时间），已远超 70 词底线且正确使用了 \textit{flyttet ind}、\textit{op på}、\textit{Heldigvis + V2 倒装} 与 \textit{fordi/at} 从句，稳拿 Modul 2 写作合格！

\vspace{0.45em}
\hrule
\vspace{0.45em}

\textbf{\color{NCDRed}【Opgave 1 范文二：A2+ 高分示范版（118 ord · 五大支柱全面秀出）】}\\
\textit{Til: Laura \quad|\quad Emne: Ny bolig}

\vspace{0.25em}
\textit{Hej Laura}

\vspace{0.2em}
\textit{Tusind tak for din besked! I sidste weekend \dabold{flyttede jeg} endelig \da{ind i} et lyst rækkehus i Roskilde, hvor jeg \dabold{har boet} i fem dage. Huset har tre værelser, et åbent køkken og en lille terrasse \da{ude} bagved.}

\vspace{0.2em}
\textit{\da{Da vi startede} flyttedagen lørdag morgen, \dabold{var jeg} lidt stresset, fordi jeg \dabold{ikke havde} nået at pakke alt ned. Heldigvis kom tre venner forbi og hjalp med at bære sofaen og alle kasserne \da{ind i} stuen.}

\vspace{0.2em}
\textit{\da{Det er især} de stille gader og gåafstanden \da{ned til} havnen, \da{der} gør mig glad for det nye kvarter. Næste lørdag den 27. maj holder jeg indflytningsfest klokken atten, og \dabold{jeg ville blive} så glad, hvis du \dabold{kunne} komme!}

\vspace{0.2em}
\textit{Mange hilsner\\
Emma}

\vspace{0.25em}
\small\textbf{高分点评（118 词）}：在稳过版基础上精准调度全册五大语法支柱：\ding{172} 过去时（\textit{flyttede}）与现在完成时（\textit{har boet}）无缝切换；\ding{173} 动静方位副词（\textit{flyttede ind i / terrasse ude bagved / ned til havnen}）；\ding{174} 前置 \textit{Da}-从句触发主句 V2 倒装（\textit{var jeg}）+ \textit{fordi} 从句中 \textit{ikke} 前置；\ding{175} 分裂强调句（\textit{Det er især..., der...}）；\ding{176} 礼貌虚拟语气邀约（\textit{jeg ville blive... hvis du kunne}）。

\vspace{0.45em}
\hrule
\vspace{0.45em}

\textbf{\color{NCDTeal}【Bonus Opgave 1B 范文：向主管 Anne 申请调班休假（101 ord）】}\\
\textit{Til: Anne (Chef) \quad|\quad Emne: Ønske om fridage i næste måned}

\vspace{0.25em}
\textit{Kære Anne}

\vspace{0.2em}
\textit{Jeg håber, at du har haft en god uge. Jeg skriver til dig, fordi jeg gerne \da{ville} spørge, om jeg \da{måtte} holde tre dages fri næste måned fra torsdag den 12. til lørdag den 14. juni.}

\vspace{0.2em}
\textit{Min søster holder nemlig en stor familiefest i Tyskland, så hele vores familie rejser derned. Jeg har allerede talt med min kollega Kasper, og han har sagt ja til at tage min torsdagsvagt for mig. \dabold{Til gengæld} vil jeg meget gerne tage et par ekstra vagter senere i juni.}

\vspace{0.2em}
\textit{\da{Ville} det måske være muligt? Det \da{ville} være en kæmpe hjælp!}

\vspace{0.2em}
\textit{Venlig hilsen\\
Emma}

\vspace{0.25em}
\small\textbf{点评（101 词）}：使用 \textit{Kære Anne} 与 \textit{Venlig hilsen} 体现半正式职场分寸；用 \textit{ville spørge, om jeg måtte} 与 \textit{Ville det måske være muligt?} 完美呈现丹麦职场最看重的礼貌过去时，先摆好同事顶班方案再主动提出 \textit{Til gengæld}（作为回报）多接班，沟通力满分。
\end{answerbox}

\subsubsection*{4.4 口语示范稿（Topic 4A \& 4B 约 130 词逐字稿 + 语音 + Opgave 2 全句表与第三人称转述）}

\begin{answerbox}
\textbf{\color{NCDNavy}【Topic 4A：Min praktik（Del 1 约 133 词脱稿陈述示范 + 核心音标）】}

\vspace{0.18em}
\textit{Jeg vil gerne fortælle om min praktik \ipa{pʰʁɑɡˈtˢiɡ}. Jeg er i praktik i et lille it-firma inde i midtbyen, hvor tolv ansatte \ipa{ˈanˌsadə} arbejder på et lyst kontor \ipa{kʰʌnˈtˢoːɐ̯ˀ}.}

\vspace{0.18em}
\textit{Mine arbejdsopgaver er både spændende og varierede. Hver morgen har vi et kort møde, og derefter tester jeg nye hjemmesider og hjælper med at svare på kundemails. Jeg møder tre dage om ugen --- tirsdag, onsdag og torsdag --- fra klokken ni til femten, i alt atten timer.}

\vspace{0.18em}
\textit{Mine kollegaer \ipa{kʰoˈleːɐ} er utrolig hjælpsomme, og min chef, som hedder Søren, er meget venlig. Det gode er, at vi kun taler dansk på kontoret, så jeg øver mit sprog hver dag. Klokken tolv holder vi en halv times frokostpause i kantinen, hvor vi spiser og drikker kaffe sammen. Normalt cykler jeg derhen på tyve minutter, men i regnvejr tager jeg bussen.}

\vspace{0.2em}
\small\textbf{[参考译文]}：我想介绍一下我的实习。我在市中心一家小 IT 公司实习，办公室很明亮，有十二名员工在那里工作。我的工作任务既有趣又多样：每天早上我们先开一个简短的晨会，之后我负责测试新网页并协助回复客户邮件。我每周出勤三天——周二、周三和周四——从上午九点到下午三点，一共十八小时。我的同事们非常乐于助人，我叫 Søren 的主管也很友善。很棒的是，我们在办公室只说丹麦语，所以我每天都能练口语。十二点我们在食堂吃半小时午饭并一起喝咖啡。平时我骑车二十分钟过去，但如果下大雨我就坐公交车。

\vspace{0.25em}
\textbf{\color{NCDRed}Del 2 考官追问与示范回答：}
\begin{enumerate}
  \item \textbf{Q:} \textit{Hvordan er det for dig, når dine kollegaer snakker hurtigt?}\\
  \textbf{A:} \textit{I starten var det lidt svært, fordi de bruger mange fagord, men jeg spørger altid pænt, om de kan gentage det lidt langsommere, og det gør de gerne.}
  \item \textbf{Q:} \textit{Hvilken arbejdsopgave synes du er mest spændende eller sværest?}\\
  \textbf{A:} \textit{Jeg synes, det er sjovest at teste nye hjemmesider. Det sværeste er at skrive kundemails helt uden stavefejl på dansk, så min kollega tjekker dem tit.}
  \item \textbf{Q:} \textit{Hvad spiser og snakker I om i frokostpausen?}\\
  \textbf{A:} \textit{Vi spiser tit varm mad eller rugbrød i kantinen, og vi snakker mest om vores weekendplaner, vejret eller gode serier.}
  \item \textbf{Q:} \textit{Kunne du tænke dig at arbejde fast der efter din praktik?}\\
  \textbf{A:} \textit{Ja, det ville være fantastisk! Jeg trives rigtig godt i teamet, og jeg håber, at der bliver en ledig stilling efter sommer.}
\end{enumerate}

\vspace{0.35em}
\hrule
\vspace{0.35em}

\textbf{\color{NCDNavy}【Topic 4B：Min fritidsaktivitet: Madklub（Del 1 约 133 词脱稿陈述示范 + 核心音标）】}

\vspace{0.18em}
\textit{I min fritid går jeg i en hyggelig madklub \ipa{ˈmæðˌkʰlub} i et medborgerhus \ipa{ˈmɛðˌbɒːɐˌhuːˀs} i Valby. Vi mødes hver anden onsdag fra klokken halv seks til klokken ni om aftenen i et stort fælleskøkken \ipa{ˈfɛləsˌkʰøɡən}. Det koster kun firs kroner pr. gang, og alle friske råvarer \ipa{ˈʁʌˌvɑːɐ} er med i prisen.}

\vspace{0.18em}
\textit{Vi er ti faste medlemmer i alle aldre. Nogle kendte jeg i forvejen, men mange var helt nye venner, da jeg begyndte for et halvt år siden. Hver gang vælger to personer et nyt land, og så laver vi mad sammen fra hele verden --- for eksempel italiensk, vietnamesisk eller dansk mad. Bagefter dækker vi et langt bord, spiser sammen og snakker i flere timer, før alle hjælper med at rydde op og vaske op. Jeg tager S-toget i ti minutter og går fem minutter fra stationen.}

\vspace{0.2em}
\small\textbf{[参考译文]}：业余时间我参加了瓦尔比市民活动中心一个温馨的美食俱乐部。我们隔周三晚上五点半到九点在一个大公共厨房聚会。每次只要八十克朗，而且所有新鲜食材都包在费用里。我们有十位不同年龄的固定成员。半年前我刚加入时，有几位本来就认识，但更多是新交的朋友。每次由两个人选一个国家，然后大家一起做来自世界各地的菜——比如意大利菜、越南菜或丹麦菜。之后我们一起摆好长桌、共进晚餐并聊上好几个小时，最后大家一起收拾洗碗。我坐十分钟轻轨，再从车站走五分钟就到。

\vspace{0.25em}
\textbf{\color{NCDRed}Del 2 考官追问与示范回答：}
\begin{enumerate}
  \item \textbf{Q:} \textit{Hvad lavede I sidste gang, du var i madklubben?}\\
  \textbf{A:} \textit{Sidste gang havde vi mexicansk tema. Jeg hjalp med at lave hjemmelavede tortillas og frisk salsa, og det smagte utrolig godt.}
  \item \textbf{Q:} \textit{Hvordan fandt du madklubben, og hvor længe har du været med?}\\
  \textbf{A:} \textit{Jeg så en plakat på det lokale bibliotek i efteråret, og jeg har været med siden november.}
  \item \textbf{Q:} \textit{Hvem bestemmer menuen, og hjælper alle med at rydde op?}\\
  \textbf{A:} \textit{Vi skiftes til at bestemme menuen to og to, og ja, alle hjælpes ad i køkkenet, så det tager kun tyve minutter at rydde op.}
  \item \textbf{Q:} \textit{Hvad betyder fællesskabet for dig og dit danske sprog?}\\
  \textbf{A:} \textit{Det betyder rigtig meget! Det er et trygt sted, hvor jeg lærer mange dagligdags ord om mad og får gode venner i nabolaget.}
\end{enumerate}
\end{answerbox}

\begin{tipbox}{口试临场救场金句（听不清考官提问 / 自我纠错 / 缓冲思考）}
\small
\begin{itemize}
  \item \textbf{请求重述或放慢}：\textit{Undskyld, jeg hørte ikke helt --- kunne du måske gentage spørgsmålet lidt langsommere?} \ipa{ˈʌnˌsgylˀ jɑ ˈhøɐ̯də eɡə ˈheːˀld}（抱歉没听清，您能稍微慢一点再重复一遍问题吗？）
  \item \textbf{争取思考时间}：\textit{Det er et rigtig godt spørgsmål... Lad mig lige tænke et sekund.}（这是个很好的问题……让我想一秒钟。）
  \item \textbf{顺滑自我纠错}：\textit{Eller det vil sige...}（或者也就是说……）/ \textit{Undskyld, jeg mener...}（抱歉，我的意思是……）——考官对主动纠错非常宽容，千万别卡壳沉默！
  \item \textbf{衔接 Opgave 2 提问}：\textit{Må jeg også spørge om, hvordan...} / \textit{Ved du tilfældigvis, hvornår han...?}（自然切换直接与间接问句，展示语言流利度。）
\end{itemize}
\end{tipbox}

\begin{tcolorbox}[enhanced, breakable, colback=NCDBgGray, colframe=NCDGray, coltitle=white, fonttitle=\bfseries, title={Mundtlig Opgave 2：Case \textit{Mikkel} 6 格提问双轨表（直接问法 + 间接问法）与 B 卡回答}, arc=2pt, boxrule=0.8pt, borderline west={3pt}{0pt}{NCDTeal}, left=8pt, right=8pt, top=6pt, bottom=6pt]
\renewcommand{\arraystretch}{1.18}
{\small
\begin{tabularx}{\linewidth}{@{} >{\bfseries\raggedright\arraybackslash}p{2.35cm} >{\raggedright\arraybackslash}p{4.05cm} >{\raggedright\arraybackslash}p{4.65cm} >{\raggedright\arraybackslash}X @{}}
\toprule
\textbf{提示词} & \textbf{直接问法（基础满分）} & \textbf{礼貌间接问法（高分展示）} & \textbf{Kursist B 回答（第三人称）} \\
\midrule
1. sted / type bolig & \textit{Hvor bor Mikkel, og hvilken type bolig bor han i?} & \textit{Kan du fortælle mig, \da{hvor} \dabold{Mikkel bor}, og hvilken slags bolig \dabold{det er}?} & \textit{Han bor i et stort rækkehus i Roskilde, ti minutter fra stationen.} \\
\midrule
2. beboere / værelser & \textit{Hvor mange bor der, og har de eget værelse?} & \textit{Ved du, \da{hvor mange} \dabold{der bor} sammen, og \da{om} \dabold{alle har} eget værelse?} & \textit{De bor fem personer sammen. Alle har eget værelse og deler køkken, stue og have.} \\
\midrule
3. regler / rengøring & \textit{Hvornår gør de rent, og spiser de sammen?} & \textit{Jeg vil gerne vide, \da{hvordan} \dabold{de gør rent}, og \da{om} \dabold{de spiser} sammen.} & \textit{De laver mad sammen tirsdag og torsdag og gør rent hver søndag efter en fast plan.} \\
\midrule
4. frivilligt arbejde / sted & \textit{Hvor arbejder han som frivillig, og hvad laver han?} & \textit{Kan du sige mig, \da{hvor} \dabold{Mikkel er frivillig}, og \da{hvad} \dabold{han laver} der?} & \textit{Han er frivillig i Røde Kors Lektiecafé på biblioteket, hvor han hjælper børn med lektier.} \\
\midrule
5. dage / tid & \textit{Hvornår arbejder han i lektiecaféen?} & \textit{Ved du, \da{hvilken dag} og \da{hvornår} \dabold{han hjælper} i lektiecaféen?} & \textit{Han er der hver mandag eftermiddag fra klokken fire til klokken seks.} \\
\midrule
6. andre hobbyer & \textit{Har han også andre fritidsinteresser?} & \textit{Jeg vil gerne høre, \da{hvad} \dabold{Mikkel ellers laver} i sin fritid.} & \textit{Ja, han løber i skoven to gange om ugen og spiller guitar i et lille band.} \\
\bottomrule
\end{tabularx}
}
\end{tcolorbox}

\begin{answerbox}
\textbf{\color{NCDTeal}【Opgave 2 第三人称总结转述示范（85 词 · 间接引语实战）】}\\
\textit{Min makker fortalte mig om Mikkel, som bor i et stort rækkehus i Roskilde tæt på stationen. Hun sagde, \da{at} han \dabold{deler} huset med fire andre personer, så de har eget værelse, men bruger køkkenet og haven sammen. De har en fast plan, hvor de laver fælles mad to gange om ugen og gør rent om søndagen. Derudover fortalte hun, \da{at} Mikkel hver mandag eftermiddag \dabold{arbejder} frivilligt i en lektiecafé på biblioteket, og \da{at} han \dabold{også løber} i skoven og spiller guitar i et band.}

\vspace{0.2em}
\small\textbf{转述亮点（85 词）}：考场上若能顺势用第三人称向考官总结搭档的回答，请注意从句连词 \textit{at} 后面保持“主语 + 动词”不倒装，并将频率副词 \textit{også} 放在动词 \textit{løber} 之前（\textit{at han også løber}），这就是从 Modul 2 迈向 Modul 3 的王牌加分动作！
\end{answerbox}
